\section{Limitações Arquiteturais}

Apesar de sua funcionalidade, a implementação atual do Patch-Hub apresenta limitações arquiteturais fundamentais que comprometem sua extensibilidade, manutenibilidade e capacidade de evolução. Essas limitações emergem de decisões de design que, embora funcionais para o escopo inicial, criam gargalos estruturais que dificultam a adição de novas funcionalidades e a manutenção do código.

\subsection{Gargalo de Mutação Centralizada}

O principal problema arquitetural reside na concentração de todo o estado da aplicação em uma única estrutura \texttt{App} que deve ser passada como \texttt{\&mut App} através de toda a cadeia de chamadas. Esta abordagem cria um \textit{gargalo de mutação centralizada} que viola princípios fundamentais de design modular.

O sistema de ownership do Rust impõe que quando uma função recebe \texttt{\&mut App}, ela adquire uma referência mutável exclusiva de toda a estrutura. Isso significa que:

\begin{itemize}
    \item Nenhum outro código pode ler ou escrever qualquer campo de \texttt{App} até que essa referência seja liberada
    \item Campos não podem ser emprestados independentemente enquanto \texttt{\&mut App} está ativa
    \item Componentes internos não podem modificar componentes irmãos sem liberar a referência primeiro
\end{itemize}

Esta restrição força a centralização de toda lógica de mutação nos métodos de \texttt{App}, impedindo que componentes gerenciem seu próprio estado de forma independente. O resultado é uma violação do princípio da responsabilidade única, onde a estrutura central deve conhecer detalhes de implementação de todos os seus componentes.

\subsection{Substituição Inadequada do Gerenciamento de Ciclo de Vida}

A arquitetura atual utiliza o padrão \texttt{Option<T>} como substituto inadequado para o gerenciamento adequado do ciclo de vida de componentes. Como a estrutura \texttt{App} sempre existe, mas as telas precisam ser inicializadas e destruídas, o código utiliza \texttt{Option<T>} para representar estados de presença:

\begin{itemize}
    \item \texttt{latest\_patchsets: Option<LatestPatchsets>}
    \item \texttt{details\_actions: Option<DetailsActions>}
    \item \texttt{edit\_config: Option<EditConfig>}
\end{itemize}

Esta abordagem introduz riscos de pânico em tempo de execução através de constantes chamadas \texttt{.as\_ref().unwrap()} e \texttt{.as\_mut().unwrap()}, substituindo a segurança em tempo de compilação oferecida pelo sistema de tipos do Rust por verificações dinâmicas propensas a falhas.

\subsection{Impossibilidade de Composição Modular}

O design atual impede a composição modular de handlers e componentes devido às restrições de borrowing. Handlers de eventos devem receber \texttt{\&mut App} inteiro mesmo quando precisam acessar apenas um subconjunto específico de dados. Isso cria uma situação onde:

\begin{itemize}
    \item Handlers não podem ser decompostos em funções auxiliares menores
    \item Não é possível passar \texttt{\&mut DetailsActions} porque \texttt{App} já possui a referência
    \item Funções auxiliares que precisam de múltiplos campos de \texttt{App} enfrentam conflitos de borrowing
\end{itemize}

Esta limitação força a criação de métodos monolíticos em \texttt{App} que conhecem detalhes de implementação de múltiplos componentes, aumentando a complexidade ciclomática e reduzindo a testabilidade.

\subsection{Atualizações de Estado Intercomponentes Problemáticas}

A arquitetura torna impossível a atualização direta de estado entre componentes, forçando todas as interações através da estrutura central \texttt{App}. Quando um componente precisa modificar o estado de outro componente irmão, ele deve passar pela estrutura \texttt{App} como intermediário, criando acoplamento forte entre componentes que deveriam ser independentes.

Por exemplo, quando \texttt{DetailsActions} precisa atualizar \texttt{BookmarkedPatchsets}, não pode manter uma referência direta a ele. Deve acessar através de \texttt{App}, tornando os componentes intimamente acoplados através do hub central.

\subsection{Concorrência e Responsividade Limitadas}

A natureza single-threaded da aplicação cria gargalos significativos durante operações de I/O, resultando em:

\begin{itemize}
    \item Interfaces não responsivas durante requisições de rede
    \item Bloqueio da UI durante operações de sistema de arquivos
    \item Lógica complexa para implementar telas de carregamento
    \item Impossibilidade de realizar operações paralelas independentes
\end{itemize}

Este design impede que a aplicação aproveite eficientemente recursos do sistema e pode levar a experiências de usuário degradadas durante operações que requerem comunicação de rede ou acesso a arquivos.

\subsection{Dificuldades de Testabilidade}

A arquitetura monolítica torna extremamente difícil a criação de testes unitários eficazes. Testes de handlers individuais requerem a construção de toda a estrutura \texttt{App}, incluindo dependências que podem não ser relevantes para o teste específico. Isso resulta em:

\begin{itemize}
    \item Testes lentos devido à necessidade de inicialização completa
    \item Testes frágeis que quebram quando componentes não relacionados são modificados
    \item Impossibilidade de isolar comportamento específico para teste
    \item Dificuldade em criar mocks e stubs para dependências externas
\end{itemize}

Esta limitação compromete significativamente a capacidade de garantir qualidade através de testes automatizados, aumentando o risco de regressões em mudanças futuras.
